\documentclass[../main.tex]{subfiles}
\begin{document}

\chapter{Compiler ILP}
\lessoninfo{
  video={Lesson 10, videos 1--20},
  textbook={H\&P \cite{hennessy2017} §3.2, Appendix H},
  keywords={instruction-level parallelism, instructions per cycle, dependence chain, tree height reduction, instruction scheduling, in-order execution, if-conversion, loop unrolling, function call inlining, code bloat, instruction count, software pipelining, trace scheduling},
}

Lessons 6--9 raised the instructions per cycle (IPC) of a processor with
hardware: register renaming, out-of-order execution with a reorder buffer, and
out-of-order load-store execution let the IPC approach the instruction-level
parallelism (ILP) of the program.  The hardware, however, can exploit only the
parallelism that the program contains and that lies within its reach.  This
lesson shows how the compiler improves both: tree height reduction shortens
the dependence chains that bound the ILP, instruction scheduling places
independent instructions next to each other, and if-conversion, loop unrolling
and function call inlining remove the branches, calls and returns that keep
them apart.\caution{\enquote{Compiler ILP} is not a second kind of ILP: the ILP
is a property of the program (Lesson 6), and the compiler changes the program.}

\section{How the compiler can raise IPC}
\label{sec:l10-limits}

The IPC of every processor is bounded by the ILP of the program it executes
(\cref{sec:l06-ipc}).\source{Lesson 10, videos 1--2; H\&P §3.2.}  Two limits
on the IPC are not determined by the processor alone, and the compiler
influences both.
\begin{description}
  \item[The ILP of the program] The ILP is the number of instructions divided
    by the length of the longest chain of RAW dependences
    (\cref{eq:l06-ilp,eq:l06-cycle}).  How long the chains are depends on how
    the compiler has translated the source program: the same computation can
    often be expressed by instructions that form shorter chains.
  \item[The parallelism a processor can find] A processor looks for
    instructions to execute only among a limited number of instructions that
    follow the oldest unexecuted one.  An out-of-order processor considers the
    instructions that fit into its reorder buffer (Lesson 8); an in-order
    processor considers only the next few and stops at the first one that must
    wait (Lesson 6).  Independent instructions that lie farther apart in
    program order than this window cannot execute in the same cycle, however
    high the ILP.  The compiler can place such instructions close together, and
    it can remove the branches, calls and returns that stand between them.
\end{description}
The compiler sees the whole program rather than a window of it, and it runs
before the program executes, so it can afford analyses that a processor could
never carry out within a clock cycle.\margin{The compiler's window is a whole
function, and after inlining the whole call tree; the processor's is its
reorder buffer (Lesson 8).}  \Cref{tab:l10-overview} lists the
techniques of this lesson.

\begin{table}[htbp]
  \centering\small
  \caption{Limits on the IPC and the compiler techniques that address them.}
  \label{tab:l10-overview}
  \begin{tabular}{@{}llc@{}}
    \toprule
    Limit & Compiler technique & Section \\
    \midrule
    long dependence chains lower the ILP    & tree height reduction
      & \ref{sec:l10-thr} \\
    independent instructions lie far apart  & instruction scheduling
      & \ref{sec:l10-scheduling} \\
    branches separate instructions          & if-conversion, trace scheduling
      & \ref{sec:l10-if-conversion}, \ref{sec:l10-other} \\
    loop branches separate iterations       & loop unrolling, software pipelining
      & \ref{sec:l10-unrolling}, \ref{sec:l10-other} \\
    calls and returns separate instructions & function call inlining
      & \ref{sec:l10-inlining} \\
    \bottomrule
  \end{tabular}
\end{table}

\section{Tree height reduction}
\label{sec:l10-thr}

A compiler translates an arithmetic expression by way of its \emph{expression
tree}, in which every operation is a node whose children are its
operands.\source{Lesson 10, videos 3--4; H\&P Appendix H.}  Each operation
becomes one instruction, and the instruction of a node reads the results of
the instructions of its children.  The longest chain of RAW dependences among
these instructions therefore equals the \emph{height} of the tree, the number
of operations on the longest path from the root to a leaf, and on the ideal
processor of Lesson 6 each level of the tree takes one cycle.\margin{One cycle
per level assumes one-cycle operations.  With a multi-cycle operation each
level costs its latency, and the height matters even more.}

A straightforward translation follows the order in which the source program
writes the operands.  For a sum of six variables held in R1 to R6, with the
result in R10, it produces
\begin{lstlisting}[language={[x86masm]Assembler}, numbers=none]
ADD R10, R1, R2      ; cycle 1
ADD R10, R10, R3     ; cycle 2
ADD R10, R10, R4     ; cycle 3
ADD R10, R10, R5     ; cycle 4
ADD R10, R10, R6     ; cycle 5
\end{lstlisting}
Every addition uses the result of the previous one.  The tree degenerates into
a chain of height~5 (\cref{fig:l10-tree-height}a): the five instructions take
five cycles on any processor, and their ILP is~1.

\begin{definition}[Tree height reduction]
  \label{def:l10-thr}
  Rewriting an expression into an equivalent one whose expression tree has
  fewer levels, by using the associativity and commutativity of its
  operations, is called \term{tree height reduction}[木の高さ削減].  The number
  of operations, and hence of instructions, does not change.
\end{definition}

\begin{example}
  \label{ex:l10-thr}
  Grouped as $((\text{R1}+\text{R2}) + (\text{R3}+\text{R4})) +
  (\text{R5}+\text{R6})$, the same sum becomes
  \begin{lstlisting}[language={[x86masm]Assembler}, numbers=none]
ADD R10, R1, R2      ; cycle 1
ADD R11, R3, R4      ; cycle 1
ADD R12, R5, R6      ; cycle 1
ADD R10, R10, R11    ; cycle 2
ADD R10, R10, R12    ; cycle 3
\end{lstlisting}
  The three additions of pairs depend on nothing and execute in cycle~1, the
  sum of the first two pairs follows in cycle~2 and the final addition in
  cycle~3 (\cref{fig:l10-tree-height}b).  There are still five instructions,
  now in three cycles: the ILP rises from 1 to $5/3 \approx 1.67$.  The price
  is two further registers, R11 and R12, for intermediate
  results.\margin{Registers are the budget for regrouping: x86-64 has 16
  architectural integer registers, A64 and RISC-V 31; beyond them the compiler
  spills to the stack.}
\end{example}

\begin{figure}[htbp]
  \centering
  % Tree height reduction of a six-operand sum: (a) left-to-right translation,
% a chain of height 5; (b) balanced tree of height 3.  Each level of
% operations is executed in one cycle of the ideal processor.
\begin{tikzpicture}[x=1cm,y=1cm, font=\footnotesize\sffamily,
  op/.style={circle, draw, fill=ln-box, minimum size=0.5cm, inner sep=0pt},
  leaf/.style={font=\scriptsize\ttfamily, inner sep=1pt, fill=white},
  edge/.style={thick, ln-accent},
  res/.style={font=\scriptsize\ttfamily, inner sep=1pt}]
  % ---- cycle rows
  \foreach \c in {1,...,5} {
    \draw[ln-rule, dotted] (-0.2,{(\c-1)*0.8}) -- (9.7,{(\c-1)*0.8});
    \node[anchor=east, text=ln-muted, font=\scriptsize\sffamily]
      at (-0.25,{(\c-1)*0.8}) {\iflangja{サイクル \c}{cycle \c}};
  }
  % ---- (a) left to right: ((((R1+R2)+R3)+R4)+R5)+R6
  \node[leaf] (r1) at (0.2,-0.8) {R1};
  \node[leaf] (r2) at (1.0,-0.8) {R2};
  \node[op] (a1) at (0.6,0.0) {$+$};
  \node[op] (a2) at (1.2,0.8) {$+$};
  \node[op] (a3) at (1.8,1.6) {$+$};
  \node[op] (a4) at (2.4,2.4) {$+$};
  \node[op] (a5) at (3.0,3.2) {$+$};
  \node[leaf] (r3) at (1.8,0.0) {R3};
  \node[leaf] (r4) at (2.4,0.8) {R4};
  \node[leaf] (r5) at (3.0,1.6) {R5};
  \node[leaf] (r6) at (3.6,2.4) {R6};
  \draw[edge] (a1) -- (r1) (a1) -- (r2)
              (a2) -- (a1) (a2) -- (r3)
              (a3) -- (a2) (a3) -- (r4)
              (a4) -- (a3) (a4) -- (r5)
              (a5) -- (a4) (a5) -- (r6);
  \node[res] (ra) at (3.0,3.85) {R10};
  \draw[edge] (a5) -- (ra);
  \node at (1.9,-1.45) {\iflangja{(a) 左から順に：高さ 5}{(a) left to right: height 5}};
  % ---- (b) balanced: ((R1+R2)+(R3+R4))+(R5+R6)
  \begin{scope}[xshift=5.6cm]
    \node[leaf] (s1) at (0.0,-0.8) {R1};
    \node[leaf] (s2) at (0.9,-0.8) {R2};
    \node[leaf] (s3) at (1.5,-0.8) {R3};
    \node[leaf] (s4) at (2.4,-0.8) {R4};
    \node[leaf] (s5) at (3.0,-0.8) {R5};
    \node[leaf] (s6) at (3.9,-0.8) {R6};
    \node[op] (p1) at (0.45,0.0) {$+$};
    \node[op] (p2) at (1.95,0.0) {$+$};
    \node[op] (p3) at (3.45,0.0) {$+$};
    \node[op] (p4) at (1.2,0.8) {$+$};
    \node[op] (p5) at (2.3,1.6) {$+$};
    \draw[edge] (p1) -- (s1) (p1) -- (s2)
                (p2) -- (s3) (p2) -- (s4)
                (p3) -- (s5) (p3) -- (s6)
                (p4) -- (p1) (p4) -- (p2)
                (p5) -- (p4) (p5) -- (p3);
    \node[res] (rb) at (2.3,2.25) {R10};
    \draw[edge] (p5) -- (rb);
    \node at (1.95,-1.45) {\iflangja{(b) 平衡木：高さ 3}{(b) balanced: height 3}};
  \end{scope}
\end{tikzpicture}

  \caption{Expression trees of the sum of R1 to R6.  (a) The left-to-right
    translation is a chain of five additions.  (b) The balanced tree of
    \cref{ex:l10-thr} has height~3.  Each level of additions executes in one
    cycle of the ideal processor.}
  \label{fig:l10-tree-height}
\end{figure}

For $n$ operands combined by one associative operation, the left-to-right
translation has height $n-1$, whereas a balanced tree has height $\lceil
\log_2 n \rceil$.  Addition and multiplication are associative, and so are
AND, OR and XOR.\caution{Floating-point addition is not associative (rounding);
compilers regroup it only when permitted, as by \texttt{-ffast-math}.}
Subtraction is not, since $(a-b)-c \ne a-(b-c)$, but an
expression with subtractions can still be regrouped when the signs are
adjusted.  Because $a - b = a + (-b)$, the terms can be grouped freely as long
as each keeps its sign, and the signs inside a group that follows a minus sign
are inverted: for example $a-b-c-d = (a-b)-(c+d)$, which reduces the height
from 3 to~2.

\begin{keypoint}
  Tree height reduction changes the dependences between instructions, not
  their number.  It shortens the chain that bounds the ILP, which pays off on
  processors that can execute several instructions per cycle.
\end{keypoint}

\section{Instruction scheduling}
\label{sec:l10-scheduling}

A processor executes independent instructions in the same cycle only if it
finds them.\source{Lesson 10, videos 5--7; H\&P §3.2.}  A compiler usually
translates a program statement by statement, so the instructions of one
statement are followed by those of the next.  When each statement forms a
chain, the independent instructions of different statements are separated by
the rest of the chain.  An out-of-order processor whose window covers both
chains finds them anyway, but an out-of-order processor with a shorter window
does not, and neither does an in-order processor, which executes the next
instructions only up to the first one that must wait.\caution{Memory holds the
compiler back as well: a load may be moved ahead of a store only if the two
addresses can be proved different.  With two pointers of the same type they
cannot, and the order stands.}  Three compiler
techniques make independent instructions easier to find: instruction
scheduling, loop unrolling (\cref{sec:l10-unrolling}) and trace scheduling
(\cref{sec:l10-other}).

\begin{definition}[Instruction scheduling]
  \label{def:l10-scheduling}
  Reordering the instructions of a program, without changing its results, so
  that instructions that can execute in the same cycle are adjacent in program
  order is called \term{instruction scheduling}[命令スケジューリング].
\end{definition}

For an in-order processor, scheduling fills the cycles in which the processor
would otherwise wait with instructions that do not depend on the waiting one.
To schedule for a processor that executes up to $w$ instructions per cycle,
the compiler can determine for every instruction the cycle in which the ideal
processor would execute it (\cref{eq:l06-cycle}) and emit the instructions in
the order of these cycles, in groups of at most $w$.  Every instruction then
still follows the instructions whose results it reads, but the reordered
program must also compute the same results as the original.  A move that
breaks a WAW or WAR dependence changes a value that some instruction uses; the
compiler then either keeps the original order or compensates:
\begin{description}
  \item[Destination register] When a moved instruction would overwrite a
    register whose old value a later instruction still reads, it writes its
    result to another, otherwise unused register instead, and the instructions
    that read the result are changed accordingly.
  \item[Address offset] When an instruction that adds a constant to a
    register, such as an index increment, moves ahead of an instruction that
    uses this register as the base of an address, the displacement of the
    address is reduced by the constant: \texttt{SW R3, 0(R1)} placed after
    \texttt{ADDI R1, R1, 4} becomes \texttt{SW R3, -4(R1)}.\tip{Moving an
    instruction past an \texttt{ADDI} of $c$ subtracts $c$ from every
    displacement on that register; moving it the other way adds $c$.}
\end{description}

\begin{example}
  \label{ex:l10-scheduling}
  With $a$ to $h$ in R1 to R8, the statements $x = (a+b)\times c - d$ and
  $y = (e-f)\times g + h$ compile into
  \begin{lstlisting}[language={[x86masm]Assembler}, numbers=none]
I1:  ADD R9, R1, R2      ; a + b
I2:  MUL R9, R9, R3      ; (a + b) * c
I3:  SUB R9, R9, R4      ; x <- (a + b) * c - d
I4:  SUB R10, R5, R6     ; e - f
I5:  MUL R10, R10, R7    ; (e - f) * g
I6:  ADD R10, R10, R8    ; y <- (e - f) * g + h
\end{lstlisting}
  The instructions form two chains, I1--I3 and I4--I6, so the ideal processor
  executes I1 and I4 in cycle~1, I2 and I5 in cycle~2, and I3 and I6 in
  cycle~3: the ILP is~2.  A two-wide in-order processor executes only I1 in
  cycle~1, because I2 must wait for it, and only I2 in cycle~2.  In cycle~3 it
  executes I3 together with I4, but then I5 and I6 again need a cycle each.
  It takes five cycles, an IPC of $6/5 = 1.2$
  (\cref{fig:l10-scheduling}a).  Scheduled in the order I1, I4, I2, I5, I3,
  I6, every pair of adjacent instructions is independent, and the same
  processor needs three cycles, an IPC of~2, equal to the ILP
  (\cref{fig:l10-scheduling}b).  The two statements use different registers,
  so the reordering breaks no WAW or WAR dependence.
\end{example}

\begin{figure}[htbp]
  \centering
  % Instruction scheduling for a two-wide in-order processor: (a) the two
% statements compiled one after the other; (b) the scheduled order.
% Arcs: RAW dependences (first statement above, second below); shaded
% groups: instructions executed in the same cycle; numbers: that cycle.
\begin{tikzpicture}[x=1cm,y=1cm, font=\footnotesize\sffamily, >={Stealth[length=1.6mm]},
  ins/.style={draw, fill=white, minimum width=0.9cm, minimum height=0.5cm, inner sep=0pt},
  dep/.style={->, thick, ln-accent},
  grp/.style={fill=ln-accent!12, rounded corners=2pt},
  cyc/.style={font=\scriptsize\sffamily, text=ln-muted}]
  % ---------------- (a) statement by statement
  \node[anchor=west] at (0.2,1.0)
    {\iflangja{(a) 文ごとに翻訳した順序：6 命令を 5 サイクルで実行}%
              {(a) statement by statement: 6 instructions in 5 cycles}};
  \foreach \xa/\xb in {1.8/1.8, 3.3/3.3, 4.8/6.3, 7.8/7.8, 9.3/9.3}
    \fill[grp] (\xa-0.6,-0.4) rectangle (\xb+0.6,0.4);
  \foreach \n/\x in {1/1.8, 2/3.3, 3/4.8, 4/6.3, 5/7.8, 6/9.3}
    \node[ins] (a\n) at (\x,0) {I\n};
  \draw[dep] (a1.north) to[bend left=40] (a2.north);
  \draw[dep] (a2.north) to[bend left=40] (a3.north);
  \draw[dep] (a4.south) to[bend right=40] (a5.south);
  \draw[dep] (a5.south) to[bend right=40] (a6.south);
  \node[cyc, anchor=east] at (1.1,-1.0) {\iflangja{サイクル}{cycle}};
  \foreach \c/\x in {1/1.8, 2/3.3, 3/4.8, 3/6.3, 4/7.8, 5/9.3}
    \node[cyc] at (\x,-1.0) {\c};
  % ---------------- (b) scheduled
  \begin{scope}[yshift=-2.8cm]
    \node[anchor=west] at (0.2,1.0)
      {\iflangja{(b) スケジューリング後：6 命令を 3 サイクルで実行}%
                {(b) scheduled: 6 instructions in 3 cycles}};
    \foreach \xa/\xb in {1.8/3.3, 4.8/6.3, 7.8/9.3}
      \fill[grp] (\xa-0.6,-0.4) rectangle (\xb+0.6,0.4);
    \foreach \n/\x in {1/1.8, 4/3.3, 2/4.8, 5/6.3, 3/7.8, 6/9.3}
      \node[ins] (b\n) at (\x,0) {I\n};
    \draw[dep] (b1.north) to[bend left=25] (b2.north);
    \draw[dep] (b2.north) to[bend left=25] (b3.north);
    \draw[dep] (b4.south) to[bend right=25] (b5.south);
    \draw[dep] (b5.south) to[bend right=25] (b6.south);
    \node[cyc, anchor=east] at (1.1,-1.0) {\iflangja{サイクル}{cycle}};
    \foreach \c/\x in {1/1.8, 1/3.3, 2/4.8, 2/6.3, 3/7.8, 3/9.3}
      \node[cyc] at (\x,-1.0) {\c};
  \end{scope}
\end{tikzpicture}

  \caption{The program of \cref{ex:l10-scheduling} on a two-wide in-order
    processor.  Arcs are RAW dependences; shaded groups execute in the same
    cycle.  (a) In the order of the statements, the chain I1--I3 holds back
    I4.  (b) After scheduling, both chains advance in every cycle.}
  \label{fig:l10-scheduling}
\end{figure}

\begin{example}
  \label{ex:l10-loop}
  The loop
  \begin{lstlisting}[language=C, numbers=none]
for (i = 0; i < n; i++)
    a[i] = a[i] + x;
\end{lstlisting}
  adds $x$ to each of the $n$ elements of an integer array.  With R1 pointing
  to the current element, R2 holding $x$, and R5 holding the address just past
  the last element, it compiles into\margin{The compiled loop has no $i$: the
  compiler keeps the address as the only induction variable and compares it
  against the end in R5 (\emph{strength reduction}).}
  \begin{lstlisting}[language={[x86masm]Assembler}, numbers=none]
Loop: LW   R3, 0(R1)      ; R3 <- a[i]
      ADD  R3, R3, R2     ; R3 <- a[i] + x
      SW   R3, 0(R1)      ; a[i] <- R3
      ADDI R1, R1, 4      ; next element
      BNE  R1, R5, Loop   ; repeat until the end
\end{lstlisting}
  A four-wide in-order processor with perfect branch prediction executes the
  LW alone, because the ADD waits for it, and then the ADD alone, because the
  SW waits for it.  The SW and the ADDI execute together, and the BNE, which
  waits for the ADDI, executes together with the LW of the next iteration.
  Every iteration takes three cycles, a CPI of $3/5 = 0.6$.  Moved up behind
  the LW, the ADDI no longer waits behind the chain, and the SW, which now
  follows the increment, needs the displacement $-4$:
  \begin{lstlisting}[language={[x86masm]Assembler}, numbers=none]
Loop: LW   R3, 0(R1)
      ADDI R1, R1, 4
      ADD  R3, R3, R2
      SW   R3, -4(R1)     ; R1 already advanced
      BNE  R1, R5, Loop
\end{lstlisting}
  This is the order of the ideal cycles, with the branch kept at the end of
  the iteration.\caution{The $-4$ is in bytes, not elements: a displacement is
  corrected by exactly the constant the base register gained, whatever the
  element size.}  The LW and the ADDI now execute together, the ADD follows,
  and the SW, the BNE and the LW and ADDI of the next iteration fill all four
  places of the next cycle.  Every iteration takes two cycles, a CPI of
  $2/5 = 0.4$.
\end{example}

Scheduling thus matters for in-order processors,\margin{A static schedule fixes
an assumed latency: the order is right while the loads hit in the cache.  On a
miss an in-order processor stalls, and only the hardware of Lessons 7--9
reacts.} for out-of-order processors
whose window is short compared with the chains in the program, and for
out-of-order processors whenever more instructions are ready than their width,
since the order decides which of them execute first (Lesson 6).

\section{Scheduling and if-conversion}
\label{sec:l10-if-conversion}

Instruction scheduling can move instructions only within a stretch of code
that is executed as a whole, and a conditional branch ends such a
stretch.\source{Lesson 10, videos 8--9; H\&P §3.2.}  An instruction that
follows the branch is executed on only one of its paths.  Moving it ahead of
the branch would execute it on the other path as well, and instructions of the
two paths cannot be placed in the same group, because only one of the paths is
executed.  The compiler therefore cannot group instructions across a branch,
even when the instructions before and after it are independent.\margin{Such a
stretch is a \emph{basic block}.  With a branch every five instructions or so in
integer code (Lesson 3), one block alone offers little to schedule.}

If-conversion\index{if-conversion} (Lesson 5) removes this boundary for a small
if-then-else.  The converted code computes the condition and both parts
unconditionally and selects the results with conditional moves, all in one
branch-free sequence, which the compiler can schedule like any other: the
instructions that compute the condition and the two parts do not depend on
the selections, so they are scheduled into the first groups according to
their own dependences, and the selections follow.  The if-converted code of \cref{ex:l05-movzn}
already has this shape.  Its first four instructions are independent of each
other and the three conditional moves depend only on them, so a four-wide
in-order processor executes all seven instructions in two cycles.

\subsection{Loops}

A loop ends every iteration with a branch back to its beginning, and this
branch cannot usefully be if-converted.  Removing it would mean executing the
body as many times as the loop could ever iterate, and each copy would need
its own predicate telling whether its iteration should have been executed; the
instructions that compute and apply these predicates would cost more than
scheduling could gain.  Loop branches are, moreover, predicted accurately, and
predicating one would execute the code that follows the loop in every
iteration (\cref{tab:l05-which}).  The loop branch therefore stays, and it
confines instruction scheduling to a single iteration: the scheduled body of
\cref{ex:l10-loop}, with two cycles per iteration on the four-wide in-order
processor, is the best order for that processor.  When the body is short, most
of the independent instructions of the loop lie in different iterations, where
the compiler cannot bring them together.  Loop unrolling addresses this
limitation.

\section{Loop unrolling}
\label{sec:l10-unrolling}

In a loop with a short body, a large fraction of the instructions executed are
loop overhead, and the work is spread over many iterations separated by the
loop branch.\source{Lesson 10, videos 10--15; H\&P §3.2.}  In the loop of
\cref{ex:l10-loop}, three of the five instructions do the work; the ADDI and
the BNE are overhead.

\begin{definition}[Loop unrolling]
  \label{def:l10-unrolling}
  Rewriting a loop so that each iteration of the new loop performs the work of
  $k$ iterations of the original loop, by repeating the body $k$ times with
  adjusted indices and advancing the index by $k$, is called \term{loop
  unrolling}[ループ展開].  The number $k$ is the \emph{unrolling factor}:
  unrolling a loop once gives $k=2$, unrolling it twice $k=3$.
\end{definition}

Unrolled with $k=2$, for an even $n$, the loop becomes
\begin{lstlisting}[language=C, numbers=none]
for (i = 0; i < n; i += 2) {
    a[i]   = a[i]   + x;
    a[i+1] = a[i+1] + x;
}
\end{lstlisting}
and compiles into
\begin{lstlisting}[language={[x86masm]Assembler}, numbers=none]
Loop: LW   R3, 0(R1)      ; a[i] <- a[i] + x
      ADD  R3, R3, R2
      SW   R3, 0(R1)
      LW   R3, 4(R1)      ; a[i+1] <- a[i+1] + x
      ADD  R3, R3, R2
      SW   R3, 4(R1)
      ADDI R1, R1, 8      ; two elements further
      BNE  R1, R5, Loop
\end{lstlisting}
Each iteration now processes two elements with eight instructions, of which
the index increment and the branch are executed once.\margin{If $n$ is not a
multiple of $k$, the compiler adds code, typically a copy of the original
loop, for the remaining elements.}  The two copies of the body are adjacent
and separated by no branch, so the compiler can schedule them together once
the second copy keeps its intermediate result in a register of its own
(\cref{sec:l10-unrolling-cpi}).

\subsection{Instruction count and ILP}
\label{sec:l10-unrolling-ilp}

Unrolling does not affect the clock period, but it changes the other two
factors of the iron law (\cref{eq:iron-law}).  The original loop executes 5
instructions per element, the unrolled loop 8 per two elements, or 4 per
element: the overhead instructions are executed half as often, and the
instruction count falls by \SI{20}{\percent}.  With a factor~$k$ the loop
executes $(3k+2)/k$ instructions per element, which approaches 3, the work
alone.\tip{Read it as $3 + 2/k$: the overhead per element halves with every
doubling of $k$, so most of the gain lies in the first few steps.}

The ILP rises as well.  On the ideal processor (\cref{eq:l06-cycle}) the
ADDIs of successive iterations form a chain of one cycle per iteration.  Every
LW of an iteration reads R1 from the ADDI of the previous iteration and
executes in the same cycle as the ADDI of its own iteration, and its ADD and
SW follow in the next two cycles; the second write to R3 in the unrolled body
is a WAW dependence that renaming removes.  The ideal processor thus executes
one iteration, or $k$ elements, per cycle, and the ILP of a long loop
approaches the number of instructions per iteration: 5 for the original loop,
8 for the unrolled one.  This relies on iterations that depend on each other
only through the index: in a loop that accumulates, such as
\texttt{s = s + a[i]}, the additions to \texttt{s} form a chain of one cycle
per element in either version.\tip{Such a reduction is unrolled with $k$
separate accumulators, added together after the loop; for floating point this
needs the reassociation permission of \cref{sec:l10-thr}.}

\subsection{Cycles per instruction}
\label{sec:l10-unrolling-cpi}

Whether the CPI falls as well depends on the processor.
\Cref{tab:l10-unrolling} compares four versions of the loop on the four-wide
in-order processor of \cref{ex:l10-loop} and on the ideal processor, both with
perfect branch prediction.\margin{A loop branch is mispredicted about once per
run, at the exit; a remainder loop adds at most one more.}  The first two versions are
those of \cref{ex:l10-loop}.  Unrolled but not scheduled, the loop still
executes one LW--ADD--SW chain after the other: the second LW joins the first
SW, the ADDI joins the second SW, and the BNE joins the LW of the next
iteration, so every iteration of eight instructions takes five cycles.  The
CPI rises from 0.6 to $5/8 = 0.625$, although the time per element falls from
3 to 2.5 cycles with the instruction count.  Scheduled in the order of the
ideal cycles, the unrolled body becomes
\begin{lstlisting}[language={[x86masm]Assembler}, numbers=none]
Loop: LW   R3, 0(R1)
      LW   R4, 4(R1)      ; second copy uses R4
      ADDI R1, R1, 8
      ADD  R3, R3, R2
      ADD  R4, R4, R2
      SW   R3, -8(R1)     ; R1 already advanced
      SW   R4, -4(R1)
      BNE  R1, R5, Loop
\end{lstlisting}
Both compensations of \cref{sec:l10-scheduling} are needed.  The second LW
now precedes the ADD and the SW of the first copy, which still need R3, so the
second copy computes its value in R4; and both stores follow the increment, so
their displacements are reduced by 8.  The processor executes both LWs and the
ADDI in the first cycle and both ADDs in the second, and the two SWs, the BNE
and the first LW of the next iteration fill the third.  From then on the
groups no longer coincide with the iterations, and in the steady state 16
instructions, two iterations, take 5 cycles: 1.25 cycles per element, a CPI of $5/16 = 0.3125$.  Compared
with the scheduled original loop, the time per element falls from 2 to 1.25
cycles, a speedup of 1.6, of which $5/4 = 1.25$ comes from the instruction
count and $0.4/0.3125 = 1.28$ from the CPI.

\begin{table}[htbp]
  \centering\small
  \caption{The loop of \cref{ex:l10-loop} before and after unrolling with
    $k=2$, each as compiled and scheduled: instructions per element, and
    cycles per element and CPI on two processors with perfect branch
    prediction.}
  \label{tab:l10-unrolling}
  \begin{tabular}{@{}lccccc@{}}
    \toprule
    & & \multicolumn{2}{c}{Four-wide in-order} & \multicolumn{2}{c}{Ideal} \\
    \cmidrule(lr){3-4}\cmidrule(l){5-6}
    Version & Instructions & Cycles & CPI & Cycles & CPI \\
    \midrule
    Original               & 5 & 3.00 & 0.600  & 1.00 & 0.200 \\
    Scheduled              & 5 & 2.00 & 0.400  & 1.00 & 0.200 \\
    Unrolled               & 4 & 2.50 & 0.625  & 0.50 & 0.125 \\
    Unrolled and scheduled & 4 & 1.25 & 0.3125 & 0.50 & 0.125 \\
    \bottomrule
  \end{tabular}
\end{table}

On the in-order processor, unrolling thus lowers the CPI through the
scheduling it makes possible: the loop branch no longer keeps the instructions
of two iterations apart, and the compiler interleaves them.  On the ideal
processor the order of the instructions does not matter.  There the CPI falls
because the ILP rises from 5 to 8, and together with the lower instruction
count the time per element halves; an out-of-order processor with a large
window profits in the same way, up to its width.  A processor that completes
one instruction per cycle has a CPI of~1 in every version and gains only the
factor 1.25 from the instruction count.  Loop unrolling therefore lowers the
instruction count on every processor, and the CPI on processors that execute
several instructions per cycle.

\subsection{Downsides of unrolling}
\label{sec:l10-bloat}

The body grows linearly with $k$, from 5 instructions to $3k+2$, whereas the
gain levels off: the instructions per element approach~3, and a real
processor cannot raise its IPC beyond its width, however high the ILP.  Going
from $k=4$ (14 instructions, 3.5 per element) to $k=8$ (26 instructions, 3.25
per element) nearly doubles the body to execute only \SI{7}{\percent} fewer
instructions, an instance of the diminishing returns of Lesson 2.

\begin{definition}[Code bloat]
  \label{def:l10-bloat}
  The growth of the program size caused by a transformation such as loop
  unrolling is called \term{code bloat}[コード肥大化].
\end{definition}

Larger code occupies more memory, and, as the lessons on caches will show, it
causes more cache misses when its instructions are fetched, which can cancel
the gain.\margin{The front end is the first limit: Skylake caches 1536 decoded
micro-operations and Zen~2 4096, and a loop unrolled past that falls back to
the decoders.  Intel optimization manual.}  Two further problems concern the
number of iterations.  When $n$ is not a multiple of $k$, the remaining
iterations need additional code.  And
unrolling presumes that the number of remaining iterations is known when an
iteration of the unrolled loop starts.  A loop whose exit is decided inside
the body, such as a search that stops at the first zero element, cannot
simply execute $k$ copies of the body per iteration, because the later copies
could perform iterations that the original loop would not; it needs additional
exit checks.  Compilers therefore unroll with small factors, and only loops
whose bodies are short.

\begin{keypoint}
  Loop unrolling executes the index increment and the loop branch once per $k$
  iterations of the original loop, which lowers the instruction count on every
  processor.  It also lowers the CPI on processors that execute several
  instructions per cycle: on in-order processors because the compiler can
  schedule the instructions of several iterations together, on processors
  that find the parallelism themselves because the ILP of loops with
  independent iterations rises.  Its costs are code that grows with $k$ and
  extra code when the number of iterations is not a multiple of $k$ or not
  known in advance.
\end{keypoint}

\section{Function call inlining}
\label{sec:l10-inlining}

A call separates the instructions of the caller from those of the called
function, much as a branch separates the two paths of an
if-then-else.\source{Lesson 10, videos 16--18; H\&P Appendix A.}  In addition,
every call executes instructions that perform none of the function's work: the
call and the return themselves, the copying of the arguments into the
registers prescribed by the calling convention and of the result out of them,
and often the saving and restoring of registers on the stack.\margin{Real
conventions pass more in registers: the System~V AMD64 ABI the first six
integer arguments, AArch64's the first eight.  The rest go on the stack.}

\begin{definition}[Function call inlining]
  \label{def:l10-inlining}
  Replacing a call by a copy of the body of the called function, adapted to
  the operands of the caller, is called \term{function call
  inlining}[インライン展開].\caution{In C and C++ \texttt{inline} is a linkage
  keyword, not an order: GCC inlines functions not declared \texttt{inline}
  from \texttt{-O2} on, and only \texttt{always\_inline} forces it.}
\end{definition}

\begin{example}
  \label{ex:l10-inlining}
  The function
  \begin{lstlisting}[language=C, numbers=none]
int sqdiff(int x, int y) { int d = x - y; return d * d; }
\end{lstlisting}
  is called as \texttt{c = sqdiff(a, b)}, with $a$, $b$ and $c$ in R8, R9 and
  R10.  The calling convention passes the arguments in R4 and R5 and returns
  the result in R2:
  \begin{lstlisting}[language={[x86masm]Assembler}, numbers=none]
        ADD  R4, R8, R0     ; x <- a
        ADD  R5, R9, R0     ; y <- b
        CALL SqDiff
        ADD  R10, R2, R0    ; c <- result
        ...
SqDiff: SUB  R2, R4, R5     ; d <- x - y
        MUL  R2, R2, R2     ; result <- d * d
        RET
\end{lstlisting}
  Inlined, the call becomes
  \begin{lstlisting}[language={[x86masm]Assembler}, numbers=none]
        SUB  R10, R8, R9    ; c <- a - b
        MUL  R10, R10, R10  ; c <- c * c
\end{lstlisting}
  Each call executes seven instructions, the inlined code two.  The argument
  and result copies also lengthen the dependence chain from $a$ and $b$ to
  $c$: on the ideal processor the call takes four cycles, the inlined code two
  (\cref{fig:l10-inlining}).
\end{example}

\begin{figure}[htbp]
  \centering
  % Function call inlining: (a) the call with its argument and result copies,
% CALL and RET (shaded: instructions that only implement the call);
% (b) the inlined body.  Circled numbers: cycle in which the ideal processor
% executes each instruction.
\begin{tikzpicture}[x=1cm,y=1cm, font=\footnotesize\sffamily, >={Stealth[length=1.6mm]},
  ins/.style={draw, fill=white, minimum width=2.5cm, minimum height=0.5cm,
    inner sep=2pt, font=\scriptsize\ttfamily},
  ovh/.style={ins, fill=ln-box, draw=ln-rule, text=ln-muted},
  cyc/.style={circle, draw=ln-accent, text=ln-accent, inner sep=0pt,
    minimum size=0.38cm, font=\scriptsize\sffamily},
  ctl/.style={->, thick, dashed, ln-muted},
  lab/.style={font=\scriptsize\sffamily, text=ln-muted}]
  % ---------------- (a) with a call
  \node[align=center] at (1.75,0.95)
    {\iflangja{(a) 呼び出し：7 命令，4 サイクル}{(a) call: 7 instructions, 4 cycles}};
  \node[lab] at (0,0.45) {\iflangja{呼び出し側}{caller}};
  \node[ovh] (c1) at (0,0.0)  {ADD  R4, R8, R0};
  \node[ovh] (c2) at (0,-0.6) {ADD  R5, R9, R0};
  \node[ovh] (c3) at (0,-1.2) {CALL SqDiff};
  \node[ovh] (c4) at (0,-2.4) {ADD  R10, R2, R0};
  \node[lab] at (3.5,-0.15) {SqDiff};
  \node[ins] (f1) at (3.5,-0.6) {SUB  R2, R4, R5};
  \node[ins] (f2) at (3.5,-1.2) {MUL  R2, R2, R2};
  \node[ovh] (f3) at (3.5,-1.8) {RET};
  \foreach \n/\c in {c1/1, c2/1, c3/1, c4/4}
    \node[cyc] at ([xshift=-0.3cm]\n.west) {\c};
  \foreach \n/\c in {f1/2, f2/3, f3/2}
    \node[cyc] at ([xshift=0.3cm]\n.east) {\c};
  \draw[ctl] (c3.east) -- ++(0.35,0) |- (f1.west);
  \draw[ctl] (f3.west) -- ++(-0.35,0) |- (c4.east);
  % ---------------- (b) inlined
  \node[align=center] at (7.3,0.95)
    {\iflangja{(b) インライン展開後：2 命令，2 サイクル}{(b) inlined: 2 instructions, 2 cycles}};
  \node[ins] (i1) at (7.3,-0.6) {SUB  R10, R8, R9};
  \node[ins] (i2) at (7.3,-1.2) {MUL  R10, R10, R10};
  \foreach \n/\c in {i1/1, i2/2}
    \node[cyc] at ([xshift=0.3cm]\n.east) {\c};
\end{tikzpicture}

  \caption{\Cref{ex:l10-inlining} before and after inlining.  Shaded
    instructions only implement the call; dashed arrows are the transfers of
    control; circled numbers give the cycle in which the ideal processor
    executes each instruction.}
  \label{fig:l10-inlining}
\end{figure}

Inlining thus lowers the instruction count, since the call, the return and
the copying of arguments and results disappear, and it shortens dependence
chains, since the body reads the caller's operands directly and delivers its
result where the caller needs it.  Moreover, the body becomes part of the
caller's code: the compiler can schedule its instructions together with those
of the caller and apply tree height reduction across the former call, and the
processor no longer has to predict the call and the return
(Lesson 4).\sidenote{Inlining is above all an \emph{enabling} transformation:
with the body in the caller, the arguments are often known constants there, so
constant propagation, dead code elimination and the techniques of this lesson
apply across the former boundary.  Compilers therefore inline early and
optimize afterwards.}

\subsection{Code bloat from inlining}
\label{sec:l10-inlining-bloat}

Inlining places a separate copy of the body at every call site, and for large
functions called from many sites the code grows considerably.

\begin{example}
  \label{ex:l10-inlining-size}
  A function with a body of 20 instructions and a return is called from 10
  sites, each of which spends 4 instructions on the call (two argument copies,
  the call and a result copy).  Without inlining the code holds $10 \times 4 +
  21 = 61$ instructions, and each call executes $4 + 21 = 25$.  Inlined at all
  sites, the function and the call sequences disappear, but every site
  receives the 20 instructions of the body: 200 instructions, more than three
  times the code, to save 5 of the 25 instructions (\SI{20}{\percent}) of each
  call.  For the function of \cref{ex:l10-inlining}, whose body has two
  instructions, the same 10 sites need $10 \times 4 + 3 = 43$ instructions
  without inlining and only $10 \times 2 = 20$ with it: the code shrinks, and
  each call executes 2 instructions instead of 7.
\end{example}

As with loop unrolling, the larger code causes more cache misses and can
cancel the gain.\margin{First-level instruction caches have stayed
small---32~KB on both Skylake and Zen~3---and every inlined copy competes for
them.}  Compilers therefore inline small functions, whose call
overhead is large compared with their body, and functions called from a single
site, whose one copy replaces the original; large functions called from many
sites remain calls.

\begin{keypoint}
  Function call inlining removes the call, the return and the passing of
  arguments and results, lowering the instruction count and shortening
  dependence chains, and it lets the compiler optimize the body together with
  the caller.  It pays off for small functions; for large functions called
  from many sites, code bloat outweighs the gain.
\end{keypoint}

\section{Other compiler techniques}
\label{sec:l10-other}

The techniques of this lesson are examples of a larger body of compiler
transformations that raise the IPC.\source{Lesson 10, videos 19--20;
\cite{lam1988}; \cite{fisher1981}.}  Two further techniques carry instruction
scheduling across the boundaries that confine it; they are only outlined here.
In \term{software pipelining}[ソフトウェアパイプライニング] the compiler treats
a loop like a pipeline: it schedules the instructions of successive iterations
so that the iterations overlap, and every cycle combines independent
instructions of several iterations while the results stay those of the
original loop \cite{lam1988}.  The second, \term{trace
scheduling}[トレーススケジューリング], carries if-conversion further: the
compiler uses a profile to select the most frequently executed path through
several branches, schedules this path as one block of code, and adds checks
that branch to compensation code, which corrects the results, when execution
leaves the path \cite{fisher1981}.  Carried to its end, scheduling by the
compiler leads to processors that leave the grouping of operations into
parallel steps entirely to the compiler: the VLIW processors of
Lesson 11.\sidenote{That VLIW supported software pipelining in hardware:
IA-64's \emph{rotating} registers rename a register by iteration, rotating
predicates switch the stages of the pipeline on and off, and a dedicated loop
branch fills and drains it, so one copy of the body does what unrolling needs
several copies for.}

\begin{keypoint}[Lesson summary]
  The IPC is limited by the ILP of the program and by the independent
  instructions a processor can find within its window; the compiler influences
  both.  Tree height reduction regroups expressions so that their instructions
  form shorter chains, raising the ILP at the same instruction count.
  Instruction scheduling makes independent instructions adjacent, changing
  destination registers and address offsets where instructions move past each
  other; it helps in-order and small-window processors.  Branches confine it,
  and if-conversion removes them from small if-then-else constructs.  Loop
  unrolling executes the index increment and the loop branch once per $k$
  iterations, lowering the instruction count; it also lowers the CPI, because
  the compiler can schedule the instructions of several iterations together
  and because the ILP of loops with independent iterations rises.  Function
  call inlining removes calls, returns and argument copying.  Both cause code
  bloat, so compilers reserve them for short loops and small functions.
  Software pipelining and trace scheduling extend scheduling across loop
  iterations and across branches.
\end{keypoint}

\end{document}
